% Zone „Das kennst du schon“ – aus bank/rekonstruktion-von-funktionsgleichungen/zone.jsonl (zusammenbau v0.1)
\einheitenkopf[zone][Kennst du schon]{Das kennst du schon}
\setcounter{aufgabe}{0}

%% TODO Titel: Ich-kann-Satz fehlt in der Bank (Platzhalter: Kettenname) – Nr. 1
%% TODO Anweisung: ein Satz über der ersten Teilaufgabe fehlt in der Bank – Nr. 1
\begin{aufgabe}{Lineare Gleichungssysteme lösen}
\begin{gleichungsraster}[2]
\gl{\text{$x + y = 7$ und $x - y = 1$ – $x$ und $y$?}} &
\gl{\text{$c = 3$, $a + b + c = 2$ und $4a + 2b + c = 9$ – $a$ und $b$?}} \\
\end{gleichungsraster}
\end{aufgabe}

%% TODO Titel: Ich-kann-Satz fehlt in der Bank (Platzhalter: Kettenname) – Nr. 2
%% TODO Anweisung: ein Satz über der ersten Teilaufgabe fehlt in der Bank – Nr. 2
\begin{aufgabe}{Lineare Funktionen}
\begin{teile}
\teil Gerade durch $A(0|3)$ und $B(2|7)$ – Steigung und $y$-Achsenabschnitt? m = \leerfeld , n = \leerfeld
\teil Gerade durch $A(-2|5)$ und $B(4|2)$ – Gleichung? y = \leerfeld
\end{teile}
\end{aufgabe}

%% TODO Titel: Ich-kann-Satz fehlt in der Bank (Platzhalter: Kettenname) – Nr. 3
%% TODO Anweisung: ein Satz über der ersten Teilaufgabe fehlt in der Bank – Nr. 3
\begin{aufgabe}{Parabelformen kennen}
\begin{teile}
\teil $f(x) = (x - 4)^2 + 2$ – Scheitelpunkt? S( \leerfeld | \leerfeld )
\teil $f(x) = -\frac{1}{2}(x + 2)^2 + 3$ – Scheitelpunkt und Öffnung? S( \leerfeld | \leerfeld ), geöffnet nach \leerfeld
\end{teile}
\end{aufgabe}

%% TODO Titel: Ich-kann-Satz fehlt in der Bank (Platzhalter: Kettenname) – Nr. 4
%% TODO Anweisung: ein Satz über der ersten Teilaufgabe fehlt in der Bank – Nr. 4
\begin{aufgabe}{Symmetrie am Term erkennen}
\begin{teile}
\teil $f(x) = 3x^4 - x^2 + 2$ – welche Symmetrie? \\ \kreuz{achsensymmetrisch zur $y$-Achse} \\ \kreuz{punktsymmetrisch zum Ursprung} \\ \kreuz{keins von beiden}
\teil $h(x) = 0{,}5x^5 - \frac{2}{3}x^3 + x$ – welche Symmetrie? \\ \kreuz{achsensymmetrisch zur $y$-Achse} \\ \kreuz{punktsymmetrisch zum Ursprung} \\ \kreuz{keins von beiden}
\end{teile}
\end{aufgabe}

%% TODO Titel: Ich-kann-Satz fehlt in der Bank (Platzhalter: Kettenname) – Nr. 5
%% TODO Anweisung: ein Satz über der ersten Teilaufgabe fehlt in der Bank – Nr. 5
\begin{aufgabe}{Ableitungsregeln anwenden}
\begin{teile}
\teil $f(x) = 4x^3 - 2x$ – $f'(x)$? f'(x) = \leerfeld
\teil $f(x) = \frac{1}{3}x^3 - \frac{3}{2}x^2 + 4$ – $f'(x)$ und $f'(-1)$? f'(x) = \leerfeld , f'(-1) = \leerfeld
\end{teile}
\end{aufgabe}

%% TODO Titel: Ich-kann-Satz fehlt in der Bank (Platzhalter: Kettenname) – Nr. 6
%% TODO Anweisung: ein Satz über der ersten Teilaufgabe fehlt in der Bank – Nr. 6
\begin{aufgabe}{Sinus- und Kosinusgraph mit Periode und Amplitude}
\begin{teile}
\teil $f(x) = 5 \cdot \mathrm{sin}(x)$ – Amplitude? \leerfeld
\teil $h(x) = 2 \cdot \mathrm{cos}\left(\frac{\pi}{2}x\right)$ – Periode und Amplitude? Periode \leerfeld , Amplitude \leerfeld
\end{teile}
\end{aufgabe}

%% TODO Titel: Ich-kann-Satz fehlt in der Bank (Platzhalter: Kettenname) – Nr. 7
%% TODO Anweisung: ein Satz über der ersten Teilaufgabe fehlt in der Bank – Nr. 7
\begin{aufgabe}{Lineare Gleichungssysteme lösen – Fehler finden}
\begin{teile}
\teil Ich finde den Fehler: Lea löst $5a - b = 9$ und $2a - b = 3$ so: \rechnung{(5a - b) - (2a - b) &= 9 - 3 \\ 3a - 2b &= 6}
\end{teile}
\end{aufgabe}

%% TODO Titel: Ich-kann-Satz fehlt in der Bank (Platzhalter: Kettenname) – Nr. 8
%% TODO Anweisung: ein Satz über der ersten Teilaufgabe fehlt in der Bank – Nr. 8
\begin{aufgabe}{Lineare Gleichungssysteme lösen – selbst rechnen}
\begin{gleichungsraster}[2]
\gl{\text{$4a - 2b = 6$ und $a - 2b = -3$ – $a$ und $b$?}} & \\
\end{gleichungsraster}
\end{aufgabe}
